% BHU PYQ -- Manually transcribed & error-corrected
% Paper: BCH-223 : Specialised Accounts-II
% Exam: B.Com. (Hons.) (Semester IV) Examination, 2015-16

\documentclass[11pt,a4paper]{article}
\usepackage[a4paper,top=2.5cm,bottom=2.5cm,left=2.8cm,right=2.8cm,headheight=14pt]{geometry}
\usepackage{amsmath,amssymb}
\usepackage[T1]{fontenc}
\usepackage[utf8]{inputenc}
\usepackage{lmodern,microtype,fancyhdr,enumitem,hyperref,lastpage,parskip}

\hypersetup{
    colorlinks=true,
    urlcolor=black,
    linkcolor=black,
    pdftitle={BCH-223: Specialised Accounts-II -- BHU B.Com. (Hons.) Sem IV 2015-16}
}

\pagestyle{fancy}
\fancyhf{}
\renewcommand{\headrulewidth}{0.4pt}
\renewcommand{\footrulewidth}{0.4pt}
\fancyhead[L]{\small\textit{Banaras Hindu University}}
\fancyhead[R]{\small\textit{BCH-223: Specialised Accounts-II}}
\fancyfoot[C]{\small Page \thepage\ of \pageref{LastPage}}

\newlist{parts}{enumerate}{2}
\setlist[parts,1]{label=(\roman*),leftmargin=*,itemsep=5pt,topsep=3pt}
\setlist[parts,2]{label=(\alph*),leftmargin=*,itemsep=3pt,topsep=2pt}
\newcommand{\pts}[1]{\hfill\textit{[#1]}}

\begin{document}

\begin{center}
  {\large\bfseries BANARAS HINDU UNIVERSITY}\\[2pt]
  {\normalsize B.Com. (Hons.) --- Semester IV Examination, 2015-16}\\[6pt]
  \rule{0.8\linewidth}{0.5pt}\\[6pt]
  {\large\bfseries Paper No. BCH-223: Specialised Accounts-II}\\[4pt]
  {\normalsize Time: 3 Hours\hfill Maximum Marks: 70}
\end{center}

\rule{\linewidth}{0.4pt}

\smallskip
\noindent\textit{Instructions: Attempt all the questions. All questions carry equal marks (14 marks each).}

\bigskip

\noindent\textbf{Question 1.} \\
Shanti Lal filed his petition of insolvency on 30th June, 2015. His books showed the following Balances:

\begin{center}
\small
\begin{tabular}{|l|r|r|}
\hline
\textbf{Particulars} & \textbf{Dr. (Rs.)} & \textbf{Cr. (Rs.)} \\
\hline
Cash in hand & 200 & \\
Fixtures and Fittings (Estimated to produce Rs. 800) & 2,400 & \\
Stock-in-Trade (Estimated to produce Rs. 12,000) & 18,000 & \\
Sundry Creditors & & 30,000 \\
Bills Payable & & 12,000 \\
Sundry Debtors: & & \\
\quad Good & 10,000 & \\
\quad Doubtful (Estimated to produce 60\%) & 20,000 & \\
\quad Bad & 20,000 & \\
Bank Overdraft & & 12,000 \\
Capital & & 16,600 \\
\hline
\textbf{Total} & \textbf{70,600} & \textbf{70,600} \\
\hline
\end{tabular}
\end{center}

Liabilities for Bills discounted Rs. 8,000, expected to rank Rs. 1,500. His household furniture etc. was valued at Rs. 2,000. He owned a house valued Rs. 8,000 and he had a mortgage on it of Rs. 6,000 at 5\% per annum. Interest was paid up to 31st Dec, 2014. Preferential creditors amounting to Rs. 450 are included in sundry creditors. Rs. 150 was due for Taxes on his house.

Prepare his Statement of Affairs and Deficiency Account.\pts{14}

\centerline{\textbf{OR}}

What are the preferential creditors according to the Presidency Towns Insolvency Act and the Provincial Insolvency Act? Prepare a proforma of the Statement of Affairs.\pts{14}

\medskip\rule{\linewidth}{0.2pt}

\pagebreak

\noindent\textbf{Question 2.} \\
As per Double Accounting System, how are the accounting entries made in regard to the following?
\begin{parts}
  \item Additions and Extension of Fixed Assets.
  \item Capital Account.
\end{parts}\pts{14}

\centerline{\textbf{OR}}

From the following particulars relating to an Electricity Company draw up:
\begin{parts}
  \item The Balance Sheet of the Company as on 31st March, 2016 under the Single Account System.
  \item The Capital Accounts and General Balance Sheet as on the same date under the Double Accounting System:
\end{parts}

\begin{center}
\small
\begin{tabular}{|l|r|}
\hline
\textbf{Particulars} & \textbf{Amount (Rs.)} \\
\hline
Nominal Capital & 3,00,000 \\
18,000 Equity Shares @ Rs. 10 & 1,80,000 \\
Debentures & 1,20,000 \\
Cash in hand and at Bank & 20,000 \\
Investments & 18,000 \\
Stock & 39,000 \\
Sundry Creditors & 9,000 \\
Sundry Debtors & 17,000 \\
Reserve Fund & 18,000 \\
\hline
\multicolumn{2}{|l|}{\textbf{Expenditure up to 31st March, 2015:}} \\
\hline
Land & 38,000 \\
Mains & 1,29,000 \\
Machinery and Plant & 53,000 \\
Building & 32,000 \\
\hline
\end{tabular}
\end{center}

The expenditure during the year ended 31st March, 2016 was Rs. 10,000; Rs. 18,000 and Rs. 8,000 on the last three items and a depreciation fund of Rs. 22,000 had been created. The balancing item of Rs. 33,000 may be taken as the profit of the company.\pts{14}

\medskip\rule{\linewidth}{0.2pt}

\pagebreak

\noindent\textbf{Question 3.} \\
Discuss the provisions for keeping accounting records in a General Insurance Company. Give the format of the Revenue Account of a General Insurance Company.\pts{14}

\centerline{\textbf{OR}}

Prepare the Revenue Account of Asha Insurance Co. Ltd. from the following information in regard to its Fire Insurance business for the year ended 31st March, 2016:

\begin{center}
\small
\begin{tabular}{|l|r|}
\hline
\textbf{Particulars} & \textbf{Amount (Rs.)} \\
\hline
Claims admitted but not paid (31.03.2016) & 42,376 \\
Commission paid & 50,000 \\
Commission on re-insurance received & 12,000 \\
Share Transfer fee & 5,000 \\
Expenses of management & 78,000 \\
Bad debts & 2,500 \\
Claims paid & 15,000 \\
Profit and Loss Appropriation Account & 10,000 \\
Premium received less re-insurance & 5,52,000 \\
Reserve for unexpired risk (01.04.2015) & 2,30,000 \\
Additional reserve (01.04.2015) & 40,000 \\
Claims outstanding (01.04.2015) & 27,000 \\
Dividend on share capital & 18,500 \\
\hline
\end{tabular}
\end{center}

The following additional information has also to be used:
\begin{parts}
  \item Premium outstanding at the end of the year Rs. 40,000.
  \item It is the policy of the company to maintain 50\% of net premium towards reserve for unexpired risks.
  \item Additional reserve at 10\% of net premium to be maintained.
\end{parts}\pts{14}

\medskip\rule{\linewidth}{0.2pt}

\noindent\textbf{Question 4.}
\begin{parts}
  \item Explain the Fundamental Principles of Government Accounting.
  \item The Bharat Gas Company re-built and re-equipped a part of their works at a cost of Rs. 7,80,000. However, the cost of this part superseded its original cost of Rs. 3,00,000. Rs. 10,000 were received from the sale of old material and Rs. 20,000 have been used in re-construction which is not included in the cost of Rs. 7,80,000. Give necessary Journal entries in the books of the company.
\end{parts}\pts{14}

\centerline{\textbf{OR}}

\begin{parts}
  \item Commercial vs. Government Accounting. Explain.
  \item Vijay filed his petition of insolvency on 31st December, 2015. The details of outstanding expenses are as follows:
  \begin{parts}
    \item Salaries of 4 Clerks @ Rs. 100 p.m. for each clerk for 3 months.
    \item Wages of one labourer for the month of August, 2015 Rs. 60.
    \item Wages for 4 labourers for the month of Sept. 2015 @ Rs. 40 per labour.
    \item Rent due to landlord for the month of Nov, 2015 Rs. 400.
    \item Amount due to the municipality for Rs. 500, to the Income Tax Department Rs. 1,000 and to the Sales Tax Department Rs. 1,000.
    \item Arrears of workmen's compensation Rs. 400.
  \end{parts}
  You are required to find out the amount of Preferential Creditors under the Provincial Insolvency Act.
\end{parts}\pts{14}

\medskip\rule{\linewidth}{0.2pt}

\pagebreak

\noindent\textbf{Question 5.}
\begin{parts}
  \item Explain the Capital Structure of a Bank.
  \item While closing the books of a Bank on 31st March, 2016 you find that in the loan ledger an Unsecured Balance of Rs. 4,00,000, in the accounts of the Bank. The debtor was bad or doubtful. Interest on the same account during the year amounted to Rs. 40,000. How would you deal with this item of Interest in 2014-15?
  
  During the year 2015-16, the bank accepts 75 paise in a rupee on account of total debt. Prepare the necessary accounts showing the ultimate effect of the transactions in 2015-16 in the accounts of the Bank.
\end{parts}\pts{14}

\centerline{\textbf{OR}}

From the following particulars prepare the Profit \& Loss Account of the Ganesh Bank Ltd. for the year ended 31st March, 2016:

\begin{center}
\footnotesize
\begin{tabular}{|l|r||l|r|}
\hline
\textbf{Particulars} & \textbf{Amount (Rs.)} & \textbf{Particulars} & \textbf{Amount (Rs.)} \\
\hline
Interest on Loans & 2,59,000 & Rent and Taxes & 18,000 \\
Interest on Fixed Deposits & 2,75,000 & Interest on Overdrafts & 54,000 \\
Rebate on Bills Discounted & 49,000 & Directors \& Auditors' Fees & 4,200 \\
Commission Charged & 8,200 & Interest on Saving Deposits & 68,000 \\
Establishment Expenses & 54,000 & Postage \& Telegrams & 1,400 \\
Discount on Bills & 95,000 & Printing \& Advertisement & 2,900 \\
Interest on Cash Credit & 2,23,000 & Sundry Charges & 1,700 \\
Interest on Current A/cs & 42,000 & & \\
\hline
\end{tabular}
\end{center}\pts{14}


\vfill
\rule{\linewidth}{0.4pt}
\begin{center}\small
\href{https://bhu-exam.vercel.app/}{Click here to visit our website}\\[4pt]
\href{https://me-aryan.vercel.app/}{Click here to visit our contributor}\\[4pt]
\href{https://quashan-me.vercel.app/}{Click here to visit our contributor}
\end{center}

\end{document}
